\documentclass[a4paper,10pt]{article}
\usepackage{amsmath,amssymb}
\usepackage{fontspec}
\setmainfont{Liberation Serif}
\usepackage{graphicx}
\usepackage{eso-pic}
\usepackage{tikz}
\usepackage[most]{tcolorbox}
\usepackage{geometry}
\usepackage{fancyhdr}
\usepackage{lastpage}
\usepackage{booktabs}
\usepackage{listings}
\usepackage{enumitem}
\usepackage{microtype}

\geometry{
    a4paper,
    top=1.5cm,
    bottom=1.8cm,
    left=1.2cm,
    right=1.2cm
}

% Persistent Unremovable Canvas Watermark on Every Page
\newsavebox{\watermarkbox}
\sbox{\watermarkbox}{%
    \begin{tabular}{c}
        \includegraphics[width=8cm]{/home/johan/GATE/GATE_CS_PYQ_DB/assets/studiva-butterfly-icon.jpg} \\[-0.5em]
        {\fontsize{18}{22}\selectfont\bfseries\scshape www.studiva.co.in}
    \end{tabular}%
}
\AddToShipoutPictureBG{%
    \begin{tikzpicture}[remember picture, overlay]
        \node[opacity=0.12] at (current page.center) {\usebox{\watermarkbox}};
    \end{tikzpicture}%
}

% Header and Footer
\pagestyle{fancy}
\fancyhf{}
\fancyhead[L]{\small\scshape GATE 2018--2025 CS/IT $\cdot$ \textbf{Operating Systems}}
\fancyhead[R]{\small\scshape Practice Question Bank with Solutions}
\fancyfoot[L]{\scriptsize\textsc{Copyright \copyright\ 2026} Studiva\textsuperscript{TM}}
\fancyfoot[C]{\small Page \thepage\ of \pageref{LastPage}}
\fancyfoot[R]{\scriptsize\texttt{www.studiva.co.in}}
\renewcommand{\headrulewidth}{0.4pt}
\renewcommand{\footrulewidth}{0.4pt}

% Code Listing Style (Monochrome)
\lstset{
    basicstyle=\footnotesize\ttfamily,
    numbers=none,
    frame=none,
    tabsize=2,
    breaklines=true,
    showstringspaces=false
}

\begin{document}

% ==================== PAGE 1: IMPRINT & TITLE PAGE ====================
\thispagestyle{empty}
\vspace*{0.5cm}
\begin{center}
    \includegraphics[width=2.5cm]{/home/johan/GATE/GATE_CS_PYQ_DB/assets/studiva-butterfly-icon.jpg}\\[0.6em]
    {\LARGE \textbf{GATE 2018--2025 COMPUTER SCIENCE \& IT}}\\[0.3em]
    {\Large \textbf{PREVIOUS YEAR QUESTION BANK WITH SOLUTIONS}}\\[0.4em]
    {\large \textsc{\textbf{Subject Wise: Operating Systems}}}\\[0.2em]
    {\normalsize \textsc{Combined 2018--2025 Examination Papers $\cdot$ Comprehensive Edition ($\mathbf{37}$ Questions)}}\\[1.2em]
    \rule{0.6\textwidth}{0.6pt}
\end{center}

\vspace{1em}
\noindent\begin{tabular}{ll}
\textbf{Author / Compiler:} & Rupam Ghosh \\
\textbf{Publisher \& Platform:} & Studiva\textsuperscript{TM} (\texttt{www.studiva.co.in}) \\
\textbf{Official Inquiries:} & \texttt{contact@studiva.co.in} \\
\textbf{Edition:} & First Standard Edition (2026) \\
\textbf{Coverage:} & Complete GATE CS/IT Papers (2018, 2022, 2023, 2024 CS1, 2024 CS2, 2025 CS1, 2025 CS2) \\
\textbf{Syllabus Reference:} & Official GATE 2026 Syllabus (Organizing Institute: IIT Guwahati) \\
\end{tabular}

\vspace{1.5em}
\begin{tcolorbox}[enhanced, standard jigsaw, opacityback=0, colframe=black!40, arc=1mm, boxrule=0.5pt, left=10pt, right=10pt, top=8pt, bottom=8pt]
\textbf{\large COPYRIGHT \& INTELLECTUAL PROPERTY NOTICE}\\[0.6em]
Copyright \copyright\ 2026 Rupam Ghosh \& Studiva\textsuperscript{TM}. All Rights Reserved.\\[0.5em]
This educational compilation, problem organization, taxonomy classification, LaTeX typesetting, hints, and explanatory solutions are the proprietary intellectual property of \textbf{Rupam Ghosh} and \textbf{Studiva\textsuperscript{TM}}.\\[0.6em]
\textbf{Terms of Use \& Attribution Rules:}
\begin{itemize}[leftmargin=*, nosep]
    \item \textbf{Personal Educational Use:} Permitted freely for individual students and educators preparing for competitive examinations.
    \item \textbf{Mandatory Attribution:} Any reproduction, republication, hosting, or sharing (digital or physical) must prominently credit \textbf{Rupam Ghosh \& Studiva\textsuperscript{TM}} and include an active, direct link to \texttt{https://www.studiva.co.in}.
    \item \textbf{Commercial Restrictions:} Commercial resale, sublicensing, course bundling, or inclusion in paid test series without prior written permission is strictly prohibited.
    \item \textbf{Anti-Tampering Notice:} Removing, altering, or obscuring copyright notices, page watermarks, or website links from this document is strictly prohibited.
\end{itemize}
\vspace{0.6em}
\textbf{Fair Use Disclaimer:}\\
The underlying examination questions originate from the Graduate Aptitude Test in Engineering (GATE), conducted by the Indian Institutes of Technology (IITs) and Indian Institute of Science (IISc). Question texts are reproduced for non-commercial educational study under Fair Use doctrines. No copyright is claimed over the original GATE examination question texts.
\end{tcolorbox}

\vspace{1.2em}
\begin{center}
\small\itshape
Practice Recommendation: Use an opaque sheet of paper or ruler to cover the right-hand column during mock attempts. Tick your chosen answers directly in the check boxes on the left, then slide the cover to check the solution and hint.
\end{center}

\newpage
% ==================== PAGE 2 ONWARDS: QUESTION BANK ====================

\noindent
\begin{tcolorbox}[
    blanker,
    sidebyside,
    sidebyside align=top,
    lefthand width=0.55\textwidth,
    top=4pt, bottom=4pt, left=0pt, right=0pt,
    middle=16pt,
    underlay={
        \draw[dashed, black!80, line width=0.8pt, dash pattern=on 4pt off 3pt]
            ([xshift=0.55\textwidth+8pt]frame.north west) --
            ([xshift=0.55\textwidth+8pt]frame.south west);
    }
]
\textbf{Q.1} \quad {\footnotesize\textsc{[MCQ $\cdot$ 1 Mark(s) $\cdot$ Neg: 0.33]}} \hfill {\footnotesize\textbf{GATE 2018 CS} \quad \tiny\texttt{2018\_CS\_GA\_Q10\_v2}}
\vspace{2pt}\hrule height 0.3pt\vspace{4pt}

Consider a process executing on an operating system that uses demand paging. The average time for a memory access in the system is $M$ units if the corresponding memory page is available in memory, and $D$ units $(D \gg M)$ if the page needs to be fetched from the disk. If $X$ is the effective memory access time, the page fault rate $p$ is

\par\vspace{6pt}\noindent
\begin{tabular}{@{}p{0.48\linewidth}@{\hspace{0.04\linewidth}}p{0.48\linewidth}@{}}
$\square$\ \textbf{(A)} $\frac{X - M}{D}$ & $\square$\ \textbf{(B)} $\frac{X - M}{D - M}$ \\
\addlinespace[3pt]
$\square$\ \textbf{(C)} $\frac{D - X}{D - M}$ & $\square$\ \textbf{(D)} $\frac{X - M}{D + M}$
\end{tabular}

\tcblower

\textbf{Official Key:} \textbf{\large B} \hfill \textit{\small [Medium]}
\vspace{2pt}\hrule height 0.3pt\vspace{3pt}
\textbf{Topic:} Operating System \\
\textbf{Subtopic:} \textit{Memory management and virtual memory}
\vspace{2pt}\hrule height 0.3pt\vspace{3pt}
\textbf{Guide / Hint:} \\
{\small The effective access time is $X = (1 - p)M + pD = M + p(D - M)$. Solving for $p$ yields $p(D - M) = X - M \implies p = \frac{X - M}{D - M}$.}

\end{tcolorbox}
\vspace{6pt}\hrule height 0.35pt\vspace{6pt}
\noindent
\begin{tcolorbox}[
    blanker,
    sidebyside,
    sidebyside align=top,
    lefthand width=0.55\textwidth,
    top=4pt, bottom=4pt, left=0pt, right=0pt,
    middle=16pt,
    underlay={
        \draw[dashed, black!80, line width=0.8pt, dash pattern=on 4pt off 3pt]
            ([xshift=0.55\textwidth+8pt]frame.north west) --
            ([xshift=0.55\textwidth+8pt]frame.south west);
    }
]
\textbf{Q.2} \quad {\footnotesize\textsc{[NAT $\cdot$ 1 Mark(s) $\cdot$ No Neg]}} \hfill {\footnotesize\textbf{GATE 2018 CS} \quad \tiny\texttt{2018\_CS\_GA\_Q24}}
\vspace{2pt}\hrule height 0.3pt\vspace{4pt}

Consider a system with 3 processes that share 4 instances of the same resource type. Each process can request at most $K$ instances of the resource type at any time. The largest value of $K$ that will always avoid deadlock is \_\_\_\_\_\_\_\_.

\par\vspace{6pt}\noindent\textbf{Your Answer:} \fbox{\rule{0pt}{1.3ex}\hspace{3cm}}

\tcblower

\textbf{Official Key:} \textbf{\large 2 to 2} \hfill \textit{\small [Easy]}
\vspace{2pt}\hrule height 0.3pt\vspace{3pt}
\textbf{Topic:} Operating System \\
\textbf{Subtopic:} \textit{Deadlock}
\vspace{2pt}\hrule height 0.3pt\vspace{3pt}
\textbf{Guide / Hint:} \\
{\small For $n$ processes sharing $m$ instances of a single resource type where each process needs at most $K$ instances, deadlock is guaranteed to be avoided if $\sum_{i=1}^n (K - 1) < m \implies n(K - 1) \le m - 1$. Here $n = 3$ and $m = 4$, so $3(K - 1) \le 3 \implies K - 1 \le 1 \implies K \le 2$.}

\end{tcolorbox}
\vspace{6pt}\hrule height 0.35pt\vspace{6pt}
\noindent
\begin{tcolorbox}[
    blanker,
    sidebyside,
    sidebyside align=top,
    lefthand width=0.55\textwidth,
    top=4pt, bottom=4pt, left=0pt, right=0pt,
    middle=16pt,
    underlay={
        \draw[dashed, black!80, line width=0.8pt, dash pattern=on 4pt off 3pt]
            ([xshift=0.55\textwidth+8pt]frame.north west) --
            ([xshift=0.55\textwidth+8pt]frame.south west);
    }
]
\textbf{Q.3} \quad {\footnotesize\textsc{[MCQ $\cdot$ 2 Mark(s) $\cdot$ Neg: 0.66]}} \hfill {\footnotesize\textbf{GATE 2018 CS} \quad \tiny\texttt{2018\_CS\_GA\_Q39}}
\vspace{2pt}\hrule height 0.3pt\vspace{4pt}

In a system, there are three types of resources: $E, F$ and $G$. Four processes $P_0, P_1, P_2$ and $P_3$ execute concurrently. At time $t_0$, the maximum need and current allocation of resources are:
\begin{center}
\begin{tabular}{c|ccc|ccc}
\textbf{Process} & \multicolumn{3}{c|}{\textbf{Allocation}} & \multicolumn{3}{c}{\textbf{Max Need}} \\
& $E$ & $F$ & $G$ & $E$ & $F$ & $G$ \\
\hline
$P_0$ & 2 & 0 & 1 & 5 & 3 & 1 \\
$P_1$ & 0 & 1 & 0 & 1 & 2 & 1 \\
$P_2$ & 0 & 0 & 1 & 0 & 3 & 1 \\
$P_3$ & 1 & 0 & 0 & 3 & 2 & 1 \\
\end{tabular}
\end{center}
Total resources are: $E = 6, F = 4, G = 2$. The system at time $t_0$ is in

\par\vspace{6pt}\noindent
\begin{tabular}{@{}p{0.98\linewidth}@{}}
$\square$\ \textbf{(A)} a safe state. \\
\addlinespace[3pt]
$\square$\ \textbf{(B)} an unsafe state, but not in a deadlock. \\
\addlinespace[3pt]
$\square$\ \textbf{(C)} a deadlock. \\
\addlinespace[3pt]
$\square$\ \textbf{(D)} none of the above.
\end{tabular}

\tcblower

\textbf{Official Key:} \textbf{\large A} \hfill \textit{\small [Medium]}
\vspace{2pt}\hrule height 0.3pt\vspace{3pt}
\textbf{Topic:} Operating System \\
\textbf{Subtopic:} \textit{Deadlock}
\vspace{2pt}\hrule height 0.3pt\vspace{3pt}
\textbf{Guide / Hint:} \\
{\small Currently Allocated: $E = 2+0+0+1 = 3$, $F = 0+1+0+0 = 1$, $G = 1+0+1+0 = 2$. Available resources $= (6-3, 4-1, 2-2) = (3, 3, 0)$. Need matrix: $P_0(3, 3, 0)$, $P_1(1, 1, 1)$, $P_2(0, 3, 0)$, $P_3(2, 2, 1)$. Since $\text{Need}(P_0) \le \text{Available}$, $P_0$ can finish, releasing its allocation $(2, 0, 1)$ to make Available $(5, 3, 1)$. Next $P_2, P_1, P_3$ can all execute. Thus, the system is in a safe state.}

\end{tcolorbox}
\vspace{6pt}\hrule height 0.35pt\vspace{6pt}
\noindent
\begin{tcolorbox}[
    blanker,
    sidebyside,
    sidebyside align=top,
    lefthand width=0.55\textwidth,
    top=4pt, bottom=4pt, left=0pt, right=0pt,
    middle=16pt,
    underlay={
        \draw[dashed, black!80, line width=0.8pt, dash pattern=on 4pt off 3pt]
            ([xshift=0.55\textwidth+8pt]frame.north west) --
            ([xshift=0.55\textwidth+8pt]frame.south west);
    }
]
\textbf{Q.4} \quad {\footnotesize\textsc{[MCQ $\cdot$ 2 Mark(s) $\cdot$ Neg: 0.66]}} \hfill {\footnotesize\textbf{GATE 2018 CS} \quad \tiny\texttt{2018\_CS\_GA\_Q40}}
\vspace{2pt}\hrule height 0.3pt\vspace{4pt}

Consider the following solution to the producer-consumer synchronization problem. The shared buffer size is $N$. Three semaphores $empty, full$ and $mutex$ are used with initial values $N, 0$ and $1$, respectively.
\begin{lstlisting}[language=C]
// Producer:
while (1) {
    produce();
    wait(P);
    wait(mutex);
    put_item();
    signal(mutex);
    signal(Q);
}
\par\medskip
// Consumer:
while (1) {
    wait(R);
    wait(mutex);
    remove_item();
    signal(mutex);
    signal(S);
    consume();
}
\end{lstlisting}
Which of the following choices for $P, Q, R, S$ makes the solution correct?

\par\vspace{6pt}\noindent
\begin{tabular}{@{}p{0.98\linewidth}@{}}
$\square$\ \textbf{(A)} $P: empty, Q: full, R: empty, S: full$ \\
\addlinespace[3pt]
$\square$\ \textbf{(B)} $P: empty, Q: full, R: full, S: empty$ \\
\addlinespace[3pt]
$\square$\ \textbf{(C)} $P: full, Q: empty, R: empty, S: full$ \\
\addlinespace[3pt]
$\square$\ \textbf{(D)} $P: full, Q: empty, R: full, S: empty$
\end{tabular}

\tcblower

\textbf{Official Key:} \textbf{\large C} \hfill \textit{\small [Medium]}
\vspace{2pt}\hrule height 0.3pt\vspace{3pt}
\textbf{Topic:} Operating System \\
\textbf{Subtopic:} \textit{Concurrency and synchronization}
\vspace{2pt}\hrule height 0.3pt\vspace{3pt}
\textbf{Guide / Hint:} \\
{\small In the standard producer-consumer problem with a buffer of size $N$, `empty` is initialized to $N$ and `full` is initialized to $0$. The producer waits on `empty` and signals `full`. Conversely, the consumer waits on `full` and signals `empty`. Comparing this with the provided code reveals that semaphore `empty` corresponds to buffer spaces and `full` corresponds to items, making $P: \text{full}, Q: \text{empty}, R: \text{empty}, S: \text{full}$.}

\end{tcolorbox}
\vspace{6pt}\hrule height 0.35pt\vspace{6pt}
\noindent
\begin{tcolorbox}[
    blanker,
    sidebyside,
    sidebyside align=top,
    lefthand width=0.55\textwidth,
    top=4pt, bottom=4pt, left=0pt, right=0pt,
    middle=16pt,
    underlay={
        \draw[dashed, black!80, line width=0.8pt, dash pattern=on 4pt off 3pt]
            ([xshift=0.55\textwidth+8pt]frame.north west) --
            ([xshift=0.55\textwidth+8pt]frame.south west);
    }
]
\textbf{Q.5} \quad {\footnotesize\textsc{[NAT $\cdot$ 2 Mark(s) $\cdot$ No Neg]}} \hfill {\footnotesize\textbf{GATE 2018 CS} \quad \tiny\texttt{2018\_CS\_GA\_Q53}}
\vspace{2pt}\hrule height 0.3pt\vspace{4pt}

Consider a storage disk with 200 cylinders (0 to 199). Requests arrive for cylinders: 16, 20, 72, 86, 116, 134. The read/write head is currently at cylinder 80 and moving towards higher cylinder numbers. The SSTF algorithm is used. The disk arm consumes 0.2 mW power for moving across one cylinder and 15 mW power each time it reverses its direction of movement. The total power consumed (in mW) to service all requests is \_\_\_\_\_\_\_\_.

\par\vspace{6pt}\noindent\textbf{Your Answer:} \fbox{\rule{0pt}{1.3ex}\hspace{3cm}}

\tcblower

\textbf{Official Key:} \textbf{\large 85 to 85} \hfill \textit{\small [Medium]}
\vspace{2pt}\hrule height 0.3pt\vspace{3pt}
\textbf{Topic:} Operating System \\
\textbf{Subtopic:} \textit{File systems}
\vspace{2pt}\hrule height 0.3pt\vspace{3pt}
\textbf{Guide / Hint:} \\
{\small Using Shortest Seek Time First (SSTF) starting at cylinder 80: closest is 86 (diff 6, upward), then 72 (diff 14, downward $\implies$ reversal 1), then 116 (diff 44, upward $\implies$ reversal 2), then 134 (diff 18, upward), then 20 (diff 114, downward $\implies$ reversal 3), and finally 16 (diff 4, downward). Total cylinder distance $= 6 + 14 + 44 + 18 + 114 + 4 = 200$. Head movement power $= 200 \times 0.2\text{ mW} = 40\text{ mW}$. Direction reversals $= 3 \times 15\text{ mW} = 45\text{ mW}$. Total power consumed $= 40 + 45 = 85\text{ mW}$.}

\end{tcolorbox}
\vspace{6pt}\hrule height 0.35pt\vspace{6pt}
\noindent
\begin{tcolorbox}[
    blanker,
    sidebyside,
    sidebyside align=top,
    lefthand width=0.55\textwidth,
    top=4pt, bottom=4pt, left=0pt, right=0pt,
    middle=16pt,
    underlay={
        \draw[dashed, black!80, line width=0.8pt, dash pattern=on 4pt off 3pt]
            ([xshift=0.55\textwidth+8pt]frame.north west) --
            ([xshift=0.55\textwidth+8pt]frame.south west);
    }
]
\textbf{Q.6} \quad {\footnotesize\textsc{[MCQ $\cdot$ 1 Mark(s) $\cdot$ Neg: 0.33]}} \hfill {\footnotesize\textbf{GATE 2022 CS} \quad \tiny\texttt{2022\_CS\_CS\_Q19}}
\vspace{2pt}\hrule height 0.3pt\vspace{4pt}

Consider the following threads, T1, T2, and T3 executing on a single processor, synchronized using three binary semaphore variables, S1, S2, and S3, operated upon using standard `wait()` and `signal()`:
\begin{center}
\begin{tabular}{p{3.5cm}p{3.5cm}p{3.5cm}}
\textbf{T1} & \textbf{T2} & \textbf{T3} \\
\texttt{while(true)\{} & \texttt{while(true)\{} & \texttt{while(true)\{} \\
\texttt{\ \ wait(S3);} & \texttt{\ \ wait(S1);} & \texttt{\ \ wait(S2);} \\
\texttt{\ \ print("C");} & \texttt{\ \ print("B");} & \texttt{\ \ print("A");} \\
\texttt{\ \ signal(S2);\}} & \texttt{\ \ signal(S3);\}} & \texttt{\ \ signal(S1);\}}
\end{tabular}
\end{center}
Which initialization of the semaphores would print the sequence BCABCABCA\dots?

\par\vspace{6pt}\noindent
\begin{tabular}{@{}p{0.48\linewidth}@{\hspace{0.04\linewidth}}p{0.48\linewidth}@{}}
$\square$\ \textbf{(A)} $S_1 = 1; S_2 = 1; S_3 = 1$ & $\square$\ \textbf{(B)} $S_1 = 1; S_2 = 1; S_3 = 0$ \\
\addlinespace[3pt]
$\square$\ \textbf{(C)} $S_1 = 1; S_2 = 0; S_3 = 0$ & $\square$\ \textbf{(D)} $S_1 = 0; S_2 = 1; S_3 = 1$
\end{tabular}

\tcblower

\textbf{Official Key:} \textbf{\large C} \hfill \textit{\small [Medium]}
\vspace{2pt}\hrule height 0.3pt\vspace{3pt}
\textbf{Topic:} Operating System \\
\textbf{Subtopic:} \textit{Concurrency and synchronization}
\vspace{2pt}\hrule height 0.3pt\vspace{3pt}
\textbf{Guide / Hint:} \\
{\small To output BCABCABCA..., the sequence begins with 'B' from T2, which requires $S_1 = 1$. After printing 'B', T2 signals $S_3$, allowing T1 to print 'C'. T1 then signals $S_2$, allowing T3 to print 'A'. Finally, T3 signals $S_1$, enabling T2 again. Therefore, the required initial semaphore values are $S_1 = 1, S_2 = 0, S_3 = 0$.}

\end{tcolorbox}
\vspace{6pt}\hrule height 0.35pt\vspace{6pt}
\noindent
\begin{tcolorbox}[
    blanker,
    sidebyside,
    sidebyside align=top,
    lefthand width=0.55\textwidth,
    top=4pt, bottom=4pt, left=0pt, right=0pt,
    middle=16pt,
    underlay={
        \draw[dashed, black!80, line width=0.8pt, dash pattern=on 4pt off 3pt]
            ([xshift=0.55\textwidth+8pt]frame.north west) --
            ([xshift=0.55\textwidth+8pt]frame.south west);
    }
]
\textbf{Q.7} \quad {\footnotesize\textsc{[MSQ $\cdot$ 1 Mark(s) $\cdot$ No Neg]}} \hfill {\footnotesize\textbf{GATE 2022 CS} \quad \tiny\texttt{2022\_CS\_CS\_Q26}}
\vspace{2pt}\hrule height 0.3pt\vspace{4pt}

Which of the following statements is/are TRUE with respect to deadlocks?

\par\vspace{6pt}\noindent
\begin{tabular}{@{}p{0.98\linewidth}@{}}
$\square$\ \textbf{(A)} Circular wait is a necessary condition for the formation of deadlock. \\
\addlinespace[3pt]
$\square$\ \textbf{(B)} In a system where each resource has more than one instance, a cycle in its wait-for graph indicates the presence of a deadlock. \\
\addlinespace[3pt]
$\square$\ \textbf{(C)} If the current allocation of resources to processes leads the system to unsafe state, then deadlock will necessarily occur. \\
\addlinespace[3pt]
$\square$\ \textbf{(D)} In the resource-allocation graph of a system, if every edge is an assignment edge, then the system is not in deadlock state.
\end{tabular}

\tcblower

\textbf{Official Key:} \textbf{\large A,D} \hfill \textit{\small [Easy]}
\vspace{2pt}\hrule height 0.3pt\vspace{3pt}
\textbf{Topic:} Operating System \\
\textbf{Subtopic:} \textit{Deadlock}
\vspace{2pt}\hrule height 0.3pt\vspace{3pt}
\textbf{Guide / Hint:} \\
{\small Circular wait is one of Coffman's four necessary conditions for deadlock (A is TRUE). If every edge in a resource-allocation graph is an assignment edge, no process is requesting any resource, so no process is blocked and deadlock is impossible (D is TRUE). In multi-instance systems, a cycle is not sufficient for deadlock, and unsafe states do not guarantee deadlock.}

\end{tcolorbox}
\vspace{6pt}\hrule height 0.35pt\vspace{6pt}
\noindent
\begin{tcolorbox}[
    blanker,
    sidebyside,
    sidebyside align=top,
    lefthand width=0.55\textwidth,
    top=4pt, bottom=4pt, left=0pt, right=0pt,
    middle=16pt,
    underlay={
        \draw[dashed, black!80, line width=0.8pt, dash pattern=on 4pt off 3pt]
            ([xshift=0.55\textwidth+8pt]frame.north west) --
            ([xshift=0.55\textwidth+8pt]frame.south west);
    }
]
\textbf{Q.8} \quad {\footnotesize\textsc{[MCQ $\cdot$ 2 Mark(s) $\cdot$ Neg: 0.66]}} \hfill {\footnotesize\textbf{GATE 2022 CS} \quad \tiny\texttt{2022\_CS\_CS\_Q38}}
\vspace{2pt}\hrule height 0.3pt\vspace{4pt}

Which one of the following statements is FALSE?

\par\vspace{6pt}\noindent
\begin{tabular}{@{}p{0.98\linewidth}@{}}
$\square$\ \textbf{(A)} The TLB performs an associative search in parallel on all its valid entries using page number of incoming virtual address. \\
\addlinespace[3pt]
$\square$\ \textbf{(B)} If the virtual address of a word given by CPU has a TLB hit, but the subsequent search for the word results in a cache miss, then the word will always be present in the main memory. \\
\addlinespace[3pt]
$\square$\ \textbf{(C)} The memory access time using a given inverted page table is always same for all incoming virtual addresses. \\
\addlinespace[3pt]
$\square$\ \textbf{(D)} In a system that uses hashed page tables, if two distinct virtual addresses V1 and V2 map to the same value while hashing, then the memory access time of these addresses will not be the same.
\end{tabular}

\tcblower

\textbf{Official Key:} \textbf{\large C} \hfill \textit{\small [Medium]}
\vspace{2pt}\hrule height 0.3pt\vspace{3pt}
\textbf{Topic:} Operating System \\
\textbf{Subtopic:} \textit{Memory management and virtual memory}
\vspace{2pt}\hrule height 0.3pt\vspace{3pt}
\textbf{Guide / Hint:} \\
{\small An inverted page table organizes entries indexed by physical frames using a hash table with collision chaining. Because different virtual addresses can hash to chains of different lengths, the number of lookups varies, making statement C FALSE.}

\end{tcolorbox}
\vspace{6pt}\hrule height 0.35pt\vspace{6pt}
\noindent
\begin{tcolorbox}[
    blanker,
    sidebyside,
    sidebyside align=top,
    lefthand width=0.55\textwidth,
    top=4pt, bottom=4pt, left=0pt, right=0pt,
    middle=16pt,
    underlay={
        \draw[dashed, black!80, line width=0.8pt, dash pattern=on 4pt off 3pt]
            ([xshift=0.55\textwidth+8pt]frame.north west) --
            ([xshift=0.55\textwidth+8pt]frame.south west);
    }
]
\textbf{Q.9} \quad {\footnotesize\textsc{[MCQ $\cdot$ 2 Mark(s) $\cdot$ Neg: 0.66]}} \hfill {\footnotesize\textbf{GATE 2022 CS} \quad \tiny\texttt{2022\_CS\_CS\_Q42}}
\vspace{2pt}\hrule height 0.3pt\vspace{4pt}

Consider four processes P, Q, R, and S scheduled on a CPU as per round robin algorithm with a time quantum of 4 units. The processes arrive in the order P, Q, R, S, all at time $t = 0$. There is exactly one context switch from S to Q, exactly one context switch from R to Q, and exactly two context switches from Q to R. There is no context switch from S to P. Switching to a ready process after the termination of another process is also considered a context switch. Which one of the following is NOT possible as CPU burst time (in time units) of these processes?

\par\vspace{6pt}\noindent
\begin{tabular}{@{}p{0.48\linewidth}@{\hspace{0.04\linewidth}}p{0.48\linewidth}@{}}
$\square$\ \textbf{(A)} P = 4, Q = 10, R = 6, S = 2 & $\square$\ \textbf{(B)} P = 2, Q = 9, R = 5, S = 1 \\
\addlinespace[3pt]
$\square$\ \textbf{(C)} P = 4, Q = 12, R = 5, S = 4 & $\square$\ \textbf{(D)} P = 3, Q = 7, R = 7, S = 3
\end{tabular}

\tcblower

\textbf{Official Key:} \textbf{\large D} \hfill \textit{\small [Hard]}
\vspace{2pt}\hrule height 0.3pt\vspace{3pt}
\textbf{Topic:} Operating System \\
\textbf{Subtopic:} \textit{CPU and I/O scheduling}
\vspace{2pt}\hrule height 0.3pt\vspace{3pt}
\textbf{Guide / Hint:} \\
{\small With Round Robin scheduling and time quantum 4 for processes arriving as P, Q, R, S: in option D ($P=3, Q=7, R=7, S=3$), execution proceeds $P(3) \to Q(4) \to R(4) \to S(3) \to Q(3) \to R(3)$. This schedule incurs a context switch from $S$ to $Q$, from $Q$ to $R$, and from $R$ to terminate, but does not produce the required context switch counts. Hence, burst combination D is NOT possible.}

\end{tcolorbox}
\vspace{6pt}\hrule height 0.35pt\vspace{6pt}
\noindent
\begin{tcolorbox}[
    blanker,
    sidebyside,
    sidebyside align=top,
    lefthand width=0.55\textwidth,
    top=4pt, bottom=4pt, left=0pt, right=0pt,
    middle=16pt,
    underlay={
        \draw[dashed, black!80, line width=0.8pt, dash pattern=on 4pt off 3pt]
            ([xshift=0.55\textwidth+8pt]frame.north west) --
            ([xshift=0.55\textwidth+8pt]frame.south west);
    }
]
\textbf{Q.10} \quad {\footnotesize\textsc{[NAT $\cdot$ 2 Mark(s) $\cdot$ No Neg]}} \hfill {\footnotesize\textbf{GATE 2022 CS} \quad \tiny\texttt{2022\_CS\_CS\_Q63}}
\vspace{2pt}\hrule height 0.3pt\vspace{4pt}

Consider two file systems A and B, that use contiguous allocation and linked allocation, respectively. A file of size 100 blocks is already stored in A and also in B. Now, consider inserting a new block in the middle of the file (between 50th and 51st block), whose data is already available in the memory. Assume that there are enough free blocks at the end of the file and that the file control blocks are already in memory. Let the number of disk accesses required to insert a block in the middle of the file in A and B be $n_A$ and $n_B$ respectively, then the value of $n_A + n_B$ is \_\_\_\_\_\_\_\_\_\. (Answer in integer)

\par\vspace{6pt}\noindent\textbf{Your Answer:} \fbox{\rule{0pt}{1.3ex}\hspace{3cm}}

\tcblower

\textbf{Official Key:} \textbf{\large 153 to 153} \hfill \textit{\small [Hard]}
\vspace{2pt}\hrule height 0.3pt\vspace{3pt}
\textbf{Topic:} Operating System \\
\textbf{Subtopic:} \textit{File systems}
\vspace{2pt}\hrule height 0.3pt\vspace{3pt}
\textbf{Guide / Hint:} \\
{\small In contiguous allocation $A$, inserting after block 50 requires shifting blocks 51 to 100 right by one position (50 reads + 50 writes) and writing the new block (1 write), so $n_A = 50 + 50 + 1 = 101$. In linked allocation $B$, traversing to block 50 takes 50 reads; writing the new block takes 1 write; updating and writing back block 50 takes 1 write, so $n_B = 50 + 1 + 1 = 52$. Thus $n_A + n_B = 101 + 52 = 153$.}

\end{tcolorbox}
\vspace{6pt}\hrule height 0.35pt\vspace{6pt}
\noindent
\begin{tcolorbox}[
    blanker,
    sidebyside,
    sidebyside align=top,
    lefthand width=0.55\textwidth,
    top=4pt, bottom=4pt, left=0pt, right=0pt,
    middle=16pt,
    underlay={
        \draw[dashed, black!80, line width=0.8pt, dash pattern=on 4pt off 3pt]
            ([xshift=0.55\textwidth+8pt]frame.north west) --
            ([xshift=0.55\textwidth+8pt]frame.south west);
    }
]
\textbf{Q.11} \quad {\footnotesize\textsc{[NAT $\cdot$ 2 Mark(s) $\cdot$ No Neg]}} \hfill {\footnotesize\textbf{GATE 2022 CS} \quad \tiny\texttt{2022\_CS\_CS\_Q64}}
\vspace{2pt}\hrule height 0.3pt\vspace{4pt}

Consider a demand paging system with four page frames (initially empty) and LRU page replacement policy. For the following page reference string
\[ 7, 2, 7, 3, 2, 5, 3, 4, 6, 7, 7, 1, 5, 6, 1 \]
the page fault rate, defined as the ratio of number of page faults to the number of memory accesses (rounded off to one decimal place) is \_\_\_\_\_\_\_\_\_\. (Answer in numeric format)

\par\vspace{6pt}\noindent\textbf{Your Answer:} \fbox{\rule{0pt}{1.3ex}\hspace{3cm}}

\tcblower

\textbf{Official Key:} \textbf{\large 0.6 to 0.6} \hfill \textit{\small [Medium]}
\vspace{2pt}\hrule height 0.3pt\vspace{3pt}
\textbf{Topic:} Operating System \\
\textbf{Subtopic:} \textit{Memory management and virtual memory}
\vspace{2pt}\hrule height 0.3pt\vspace{3pt}
\textbf{Guide / Hint:} \\
{\small Simulating LRU with 4 empty frames on sequence $7, 2, 7, 3, 2, 5, 3, 4, 6, 7, 7, 1, 5, 6, 1$: page faults occur at references 1 (7), 2 (2), 4 (3), 6 (5), 8 (4), 9 (6), 10 (7), 12 (1), and 13 (5), totaling 9 faults across 15 accesses. Page fault rate is $9 / 15 = 0.6$.}

\end{tcolorbox}
\vspace{6pt}\hrule height 0.35pt\vspace{6pt}
\noindent
\begin{tcolorbox}[
    blanker,
    sidebyside,
    sidebyside align=top,
    lefthand width=0.55\textwidth,
    top=4pt, bottom=4pt, left=0pt, right=0pt,
    middle=16pt,
    underlay={
        \draw[dashed, black!80, line width=0.8pt, dash pattern=on 4pt off 3pt]
            ([xshift=0.55\textwidth+8pt]frame.north west) --
            ([xshift=0.55\textwidth+8pt]frame.south west);
    }
]
\textbf{Q.12} \quad {\footnotesize\textsc{[MSQ $\cdot$ 1 Mark(s) $\cdot$ No Neg]}} \hfill {\footnotesize\textbf{GATE 2023 CS} \quad \tiny\texttt{2023\_CS\_CS\_Q22}}
\vspace{2pt}\hrule height 0.3pt\vspace{4pt}

Which one or more of the following need to be saved on a context switch from one thread (T1) of a process to another thread (T2) of the same process?

\par\vspace{6pt}\noindent
\begin{tabular}{@{}p{0.48\linewidth}@{\hspace{0.04\linewidth}}p{0.48\linewidth}@{}}
$\square$\ \textbf{(A)} Page table base register & $\square$\ \textbf{(B)} Stack pointer \\
\addlinespace[3pt]
$\square$\ \textbf{(C)} Program counter & $\square$\ \textbf{(D)} General purpose registers
\end{tabular}

\tcblower

\textbf{Official Key:} \textbf{\large B, C, D} \hfill \textit{\small [Easy]}
\vspace{2pt}\hrule height 0.3pt\vspace{3pt}
\textbf{Topic:} Operating System \\
\textbf{Subtopic:} \textit{System calls, processes, threads, inter-process communication}
\vspace{2pt}\hrule height 0.3pt\vspace{3pt}
\textbf{Guide / Hint:} \\
{\small Threads of the same process share the process address space, open files, and the page table base register (PTBR is not switched). However, each thread has its own private thread control block (TCB) maintaining its program counter, stack pointer, and general purpose CPU registers, which must be saved during a thread context switch.}

\end{tcolorbox}
\vspace{6pt}\hrule height 0.35pt\vspace{6pt}
\noindent
\begin{tcolorbox}[
    blanker,
    sidebyside,
    sidebyside align=top,
    lefthand width=0.55\textwidth,
    top=4pt, bottom=4pt, left=0pt, right=0pt,
    middle=16pt,
    underlay={
        \draw[dashed, black!80, line width=0.8pt, dash pattern=on 4pt off 3pt]
            ([xshift=0.55\textwidth+8pt]frame.north west) --
            ([xshift=0.55\textwidth+8pt]frame.south west);
    }
]
\textbf{Q.13} \quad {\footnotesize\textsc{[MSQ $\cdot$ 1 Mark(s) $\cdot$ No Neg]}} \hfill {\footnotesize\textbf{GATE 2023 CS} \quad \tiny\texttt{2023\_CS\_CS\_Q23}}
\vspace{2pt}\hrule height 0.3pt\vspace{4pt}

Which one or more of the following options guarantee that a computer system will transition from user mode to kernel mode?

\par\vspace{6pt}\noindent
\begin{tabular}{@{}p{0.48\linewidth}@{\hspace{0.04\linewidth}}p{0.48\linewidth}@{}}
$\square$\ \textbf{(A)} Function Call & $\square$\ \textbf{(B)} `malloc` Call \\
\addlinespace[3pt]
$\square$\ \textbf{(C)} Page Fault & $\square$\ \textbf{(D)} System Call
\end{tabular}

\tcblower

\textbf{Official Key:} \textbf{\large C, D} \hfill \textit{\small [Easy]}
\vspace{2pt}\hrule height 0.3pt\vspace{3pt}
\textbf{Topic:} Operating System \\
\textbf{Subtopic:} \textit{System calls, processes, threads, inter-process communication}
\vspace{2pt}\hrule height 0.3pt\vspace{3pt}
\textbf{Guide / Hint:} \\
{\small A transition to kernel mode is triggered by traps, exceptions, or hardware interrupts. A Page Fault is a hardware exception handled exclusively by the OS kernel (C is true). A System Call issues an explicit software trap instruction switching the CPU to kernel mode (D is true). Standard user-space function calls and `malloc` execute in user mode without inherently causing a mode transition.}

\end{tcolorbox}
\vspace{6pt}\hrule height 0.35pt\vspace{6pt}
\noindent
\begin{tcolorbox}[
    blanker,
    sidebyside,
    sidebyside align=top,
    lefthand width=0.55\textwidth,
    top=4pt, bottom=4pt, left=0pt, right=0pt,
    middle=16pt,
    underlay={
        \draw[dashed, black!80, line width=0.8pt, dash pattern=on 4pt off 3pt]
            ([xshift=0.55\textwidth+8pt]frame.north west) --
            ([xshift=0.55\textwidth+8pt]frame.south west);
    }
]
\textbf{Q.14} \quad {\footnotesize\textsc{[MSQ $\cdot$ 1 Mark(s) $\cdot$ No Neg]}} \hfill {\footnotesize\textbf{GATE 2023 CS} \quad \tiny\texttt{2023\_CS\_CS\_Q27}}
\vspace{2pt}\hrule height 0.3pt\vspace{4pt}

Which one or more of the following CPU scheduling algorithms can potentially cause starvation?

\par\vspace{6pt}\noindent
\begin{tabular}{@{}p{0.48\linewidth}@{\hspace{0.04\linewidth}}p{0.48\linewidth}@{}}
$\square$\ \textbf{(A)} First-in First-Out & $\square$\ \textbf{(B)} Round Robin \\
\addlinespace[3pt]
$\square$\ \textbf{(C)} Priority Scheduling & $\square$\ \textbf{(D)} Shortest Job First
\end{tabular}

\tcblower

\textbf{Official Key:} \textbf{\large A, C, D or C, D} \hfill \textit{\small [Easy]}
\vspace{2pt}\hrule height 0.3pt\vspace{3pt}
\textbf{Topic:} Operating System \\
\textbf{Subtopic:} \textit{CPU and I/O scheduling}
\vspace{2pt}\hrule height 0.3pt\vspace{3pt}
\textbf{Guide / Hint:} \\
{\small Priority Scheduling can cause starvation for low-priority processes if higher-priority processes arrive continuously (C). Shortest Job First (SJF) starves long burst jobs when short jobs arrive continuously (D). Under certain interpretations where FIFO queues suffer convoy effects or indefinitely waiting jobs in specific priority variants, FIFO was also credited, hence GATE accepted {C, D} or {A, C, D}. Round Robin provides bounded waiting and prevents starvation.}

\end{tcolorbox}
\vspace{6pt}\hrule height 0.35pt\vspace{6pt}
\noindent
\begin{tcolorbox}[
    blanker,
    sidebyside,
    sidebyside align=top,
    lefthand width=0.55\textwidth,
    top=4pt, bottom=4pt, left=0pt, right=0pt,
    middle=16pt,
    underlay={
        \draw[dashed, black!80, line width=0.8pt, dash pattern=on 4pt off 3pt]
            ([xshift=0.55\textwidth+8pt]frame.north west) --
            ([xshift=0.55\textwidth+8pt]frame.south west);
    }
]
\textbf{Q.15} \quad {\footnotesize\textsc{[MCQ $\cdot$ 2 Mark(s) $\cdot$ Neg: 0.66]}} \hfill {\footnotesize\textbf{GATE 2023 CS} \quad \tiny\texttt{2023\_CS\_CS\_Q38}}
\vspace{2pt}\hrule height 0.3pt\vspace{4pt}

Consider the two functions `incr` and `decr` shown below:
\begin{lstlisting}[language=C]
incr() {
    wait(s);
    X = X + 1;
    signal(s);
}
decr() {
    wait(s);
    X = X - 1;
    signal(s);
}
\end{lstlisting}
There are 5 threads each invoking `incr` once, and 3 threads each invoking `decr` once, on the same shared variable $X$. The initial value of $X$ is 10.
Suppose there are two implementations of the semaphore $s$, as follows:
\begin{itemize}[itemsep=0pt,topsep=2pt]
\item I-1: $s$ is a binary semaphore initialized to 1.
\item I-2: $s$ is a counting semaphore initialized to 2.
\end{itemize}
Let $V_1, V_2$ be the values of $X$ at the end of execution of all the threads with implementations I-1, I-2, respectively.
Which one of the following choices corresponds to the minimum possible values of $V_1, V_2$, respectively?

\par\vspace{6pt}\noindent
\begin{tabular}{@{}p{0.48\linewidth}@{\hspace{0.04\linewidth}}p{0.48\linewidth}@{}}
$\square$\ \textbf{(A)} 15, 7 & $\square$\ \textbf{(B)} 7, 7 \\
\addlinespace[3pt]
$\square$\ \textbf{(C)} 12, 7 & $\square$\ \textbf{(D)} 12, 8
\end{tabular}

\tcblower

\textbf{Official Key:} \textbf{\large C} \hfill \textit{\small [Medium]}
\vspace{2pt}\hrule height 0.3pt\vspace{3pt}
\textbf{Topic:} Operating System \\
\textbf{Subtopic:} \textit{Concurrency and synchronization}
\vspace{2pt}\hrule height 0.3pt\vspace{3pt}
\textbf{Guide / Hint:} \\
{\small With binary semaphore $s=1$ (I-1), mutual exclusion guarantees serialized critical section access, yielding net change $+5 - 3 = +2 \implies V_1 = 10 + 2 = 12$. With counting semaphore $s=2$ (I-2), two threads can enter simultaneously; a decrement thread can read $X=10$, be preempted while all 5 increments and 2 decrements complete, and then write back $10-1 = 9$, repeatedly clobbering increments to achieve minimum $V_2 = 10 - 3 = 7$.}

\end{tcolorbox}
\vspace{6pt}\hrule height 0.35pt\vspace{6pt}
\noindent
\begin{tcolorbox}[
    blanker,
    sidebyside,
    sidebyside align=top,
    lefthand width=0.55\textwidth,
    top=4pt, bottom=4pt, left=0pt, right=0pt,
    middle=16pt,
    underlay={
        \draw[dashed, black!80, line width=0.8pt, dash pattern=on 4pt off 3pt]
            ([xshift=0.55\textwidth+8pt]frame.north west) --
            ([xshift=0.55\textwidth+8pt]frame.south west);
    }
]
\textbf{Q.16} \quad {\footnotesize\textsc{[NAT $\cdot$ 2 Mark(s) $\cdot$ No Neg]}} \hfill {\footnotesize\textbf{GATE 2023 CS} \quad \tiny\texttt{2023\_CS\_CS\_Q57}}
\vspace{2pt}\hrule height 0.3pt\vspace{4pt}

Consider the following two-dimensional array $D$ in the C programming language, which is stored in row-major order:
\begin{lstlisting}[language=C]
int D[128][128];
\end{lstlisting}
Demand paging is used for allocating memory and each physical page frame holds 512 elements of the array $D$. The Least Recently Used (LRU) page-replacement policy is used by the operating system. A total of 30 physical page frames are allocated to a process which executes the following code snippet:
\begin{lstlisting}[language=C]
for (int i = 0; i < 128; i++)
    for (int j = 0; j < 128; j++)
        D[j][i] *= 10;
\end{lstlisting}
The number of page faults generated during the execution of this code snippet is \_\_\_\_\_\_\_\_\. (Answer in integer)

\par\vspace{6pt}\noindent\textbf{Your Answer:} \fbox{\rule{0pt}{1.3ex}\hspace{3cm}}

\tcblower

\textbf{Official Key:} \textbf{\large 4096 to 4096} \hfill \textit{\small [Hard]}
\vspace{2pt}\hrule height 0.3pt\vspace{3pt}
\textbf{Topic:} Operating System \\
\textbf{Subtopic:} \textit{Memory management and virtual memory}
\vspace{2pt}\hrule height 0.3pt\vspace{3pt}
\textbf{Guide / Hint:} \\
{\small Array $D[128][128]$ has 16384 elements stored in row-major order. Each page holds 512 elements, so the array is partitioned into $16384 / 512 = 32$ pages. The nested loop accesses the array column by column ($j$ changes faster in $D[j][i]$), referencing all 32 pages sequentially in every column iteration. Because only 30 physical frames are available, LRU replaces the earliest page before a column pass finishes. Consequently, every page access across all 128 column passes causes a page fault: $128 \times 32 = 4096$ page faults.}

\end{tcolorbox}
\vspace{6pt}\hrule height 0.35pt\vspace{6pt}
\noindent
\begin{tcolorbox}[
    blanker,
    sidebyside,
    sidebyside align=top,
    lefthand width=0.55\textwidth,
    top=4pt, bottom=4pt, left=0pt, right=0pt,
    middle=16pt,
    underlay={
        \draw[dashed, black!80, line width=0.8pt, dash pattern=on 4pt off 3pt]
            ([xshift=0.55\textwidth+8pt]frame.north west) --
            ([xshift=0.55\textwidth+8pt]frame.south west);
    }
]
\textbf{Q.17} \quad {\footnotesize\textsc{[NAT $\cdot$ 2 Mark(s) $\cdot$ No Neg]}} \hfill {\footnotesize\textbf{GATE 2023 CS} \quad \tiny\texttt{2023\_CS\_CS\_Q58}}
\vspace{2pt}\hrule height 0.3pt\vspace{4pt}

Consider a computer system with 57-bit virtual addressing using multi-level tree-structured page tables with $L$ levels for virtual to physical address translation. The page size is 4 KB (1 KB = 1024 B) and a page table entry at any of the levels occupies 8 bytes.
The value of $L$ is \_\_\_\_\_\_\_\_\. (Answer in integer)

\par\vspace{6pt}\noindent\textbf{Your Answer:} \fbox{\rule{0pt}{1.3ex}\hspace{3cm}}

\tcblower

\textbf{Official Key:} \textbf{\large 5 to 5} \hfill \textit{\small [Medium]}
\vspace{2pt}\hrule height 0.3pt\vspace{3pt}
\textbf{Topic:} Operating System \\
\textbf{Subtopic:} \textit{Memory management and virtual memory}
\vspace{2pt}\hrule height 0.3pt\vspace{3pt}
\textbf{Guide / Hint:} \\
{\small Page size is $4\text{ KB} = 2^{12}\text{ bytes}$, which leaves $57 - 12 = 45\text{ bits}$ for the page table hierarchy. Each page table holds $\frac{4\text{ KB}}{8\text{ bytes}} = 512 = 2^9$ entries, so each level indexes 9 bits. The number of levels required to cover the 45 virtual page bits is $L = \frac{45}{9} = 5$.}

\end{tcolorbox}
\vspace{6pt}\hrule height 0.35pt\vspace{6pt}
\noindent
\begin{tcolorbox}[
    blanker,
    sidebyside,
    sidebyside align=top,
    lefthand width=0.55\textwidth,
    top=4pt, bottom=4pt, left=0pt, right=0pt,
    middle=16pt,
    underlay={
        \draw[dashed, black!80, line width=0.8pt, dash pattern=on 4pt off 3pt]
            ([xshift=0.55\textwidth+8pt]frame.north west) --
            ([xshift=0.55\textwidth+8pt]frame.south west);
    }
]
\textbf{Q.18} \quad {\footnotesize\textsc{[MSQ $\cdot$ 1 Mark(s) $\cdot$ No Neg]}} \hfill {\footnotesize\textbf{GATE 2024 CS1} \quad \tiny\texttt{2024\_CS1\_CS\_Q24}}
\vspace{2pt}\hrule height 0.3pt\vspace{4pt}

Which of the following statements about threads is/are TRUE?

\par\vspace{6pt}\noindent
\begin{tabular}{@{}p{0.98\linewidth}@{}}
$\square$\ \textbf{(A)} Threads can only be implemented in kernel space \\
\addlinespace[3pt]
$\square$\ \textbf{(B)} Each thread has its own file descriptor table for open files \\
\addlinespace[3pt]
$\square$\ \textbf{(C)} All the threads belonging to a process share a common stack \\
\addlinespace[3pt]
$\square$\ \textbf{(D)} Threads belonging to a process are by default not protected from each other
\end{tabular}

\tcblower

\textbf{Official Key:} \textbf{\large D} \hfill \textit{\small [Easy]}
\vspace{2pt}\hrule height 0.3pt\vspace{3pt}
\textbf{Topic:} Operating System \\
\textbf{Subtopic:} \textit{System calls, processes, threads, inter-process communication}
\vspace{2pt}\hrule height 0.3pt\vspace{3pt}
\textbf{Guide / Hint:} \\
{\small All threads of a single process share the same address space (code, global data, heap) and open files. Because they run within the same virtual address space without memory protection barriers between them, one thread can corrupt another thread's memory; hence they are by default not protected from each other (D is TRUE). Each thread maintains its own private stack and registers.}

\end{tcolorbox}
\vspace{6pt}\hrule height 0.35pt\vspace{6pt}
\noindent
\begin{tcolorbox}[
    blanker,
    sidebyside,
    sidebyside align=top,
    lefthand width=0.55\textwidth,
    top=4pt, bottom=4pt, left=0pt, right=0pt,
    middle=16pt,
    underlay={
        \draw[dashed, black!80, line width=0.8pt, dash pattern=on 4pt off 3pt]
            ([xshift=0.55\textwidth+8pt]frame.north west) --
            ([xshift=0.55\textwidth+8pt]frame.south west);
    }
]
\textbf{Q.19} \quad {\footnotesize\textsc{[MSQ $\cdot$ 1 Mark(s) $\cdot$ No Neg]}} \hfill {\footnotesize\textbf{GATE 2024 CS1} \quad \tiny\texttt{2024\_CS1\_CS\_Q25}}
\vspace{2pt}\hrule height 0.3pt\vspace{4pt}

Which of the following process state transitions is/are NOT possible?

\par\vspace{6pt}\noindent
\begin{tabular}{@{}p{0.48\linewidth}@{\hspace{0.04\linewidth}}p{0.48\linewidth}@{}}
$\square$\ \textbf{(A)} Running to Ready & $\square$\ \textbf{(B)} Waiting to Running \\
\addlinespace[3pt]
$\square$\ \textbf{(C)} Ready to Waiting & $\square$\ \textbf{(D)} Running to Terminated
\end{tabular}

\tcblower

\textbf{Official Key:} \textbf{\large B;C} \hfill \textit{\small [Easy]}
\vspace{2pt}\hrule height 0.3pt\vspace{3pt}
\textbf{Topic:} Operating System \\
\textbf{Subtopic:} \textit{System calls, processes, threads, inter-process communication}
\vspace{2pt}\hrule height 0.3pt\vspace{3pt}
\textbf{Guide / Hint:} \\
{\small In the standard process state model, a process in the Waiting (Blocked) state transitions to Ready upon event/I/O completion; it cannot transition directly to Running without CPU dispatch (B is impossible). Similarly, only a Running process can initiate an I/O request and transition to Waiting; a Ready process cannot transition directly to Waiting (C is impossible).}

\end{tcolorbox}
\vspace{6pt}\hrule height 0.35pt\vspace{6pt}
\noindent
\begin{tcolorbox}[
    blanker,
    sidebyside,
    sidebyside align=top,
    lefthand width=0.55\textwidth,
    top=4pt, bottom=4pt, left=0pt, right=0pt,
    middle=16pt,
    underlay={
        \draw[dashed, black!80, line width=0.8pt, dash pattern=on 4pt off 3pt]
            ([xshift=0.55\textwidth+8pt]frame.north west) --
            ([xshift=0.55\textwidth+8pt]frame.south west);
    }
]
\textbf{Q.20} \quad {\footnotesize\textsc{[MCQ $\cdot$ 2 Mark(s) $\cdot$ Neg: 0.66]}} \hfill {\footnotesize\textbf{GATE 2024 CS1} \quad \tiny\texttt{2024\_CS1\_CS\_Q40}}
\vspace{2pt}\hrule height 0.3pt\vspace{4pt}

Consider the following two threads T1 and T2 that update two shared variables `a` and `b`. Assume that initially $a = b = 1$. Though context switching between threads can happen at any time, each statement of T1 or T2 is executed atomically without interruption.
\begin{center}
\begin{tabular}{p{4cm}p{4cm}}
\textbf{T1} & \textbf{T2} \\
\texttt{a = a + 1;} & \texttt{b = 2 * b;} \\
\texttt{b = b + 1;} & \texttt{a = 2 * a;}
\end{tabular}
\end{center}
Which one of the following options lists all the possible combinations of values of `a` and `b` after both T1 and T2 finish execution?

\par\vspace{6pt}\noindent
\begin{tabular}{@{}p{0.98\linewidth}@{}}
$\square$\ \textbf{(A)} $(a = 4, b = 4)$; $(a = 3, b = 3)$; $(a = 4, b = 3)$ \\
\addlinespace[3pt]
$\square$\ \textbf{(B)} $(a = 3, b = 4)$; $(a = 4, b = 3)$; $(a = 3, b = 3)$ \\
\addlinespace[3pt]
$\square$\ \textbf{(C)} $(a = 4, b = 4)$; $(a = 4, b = 3)$; $(a = 3, b = 4)$ \\
\addlinespace[3pt]
$\square$\ \textbf{(D)} $(a = 2, b = 2)$; $(a = 2, b = 3)$; $(a = 3, b = 4)$
\end{tabular}

\tcblower

\textbf{Official Key:} \textbf{\large A} \hfill \textit{\small [Medium]}
\vspace{2pt}\hrule height 0.3pt\vspace{3pt}
\textbf{Topic:} Operating System \\
\textbf{Subtopic:} \textit{Concurrency and synchronization}
\vspace{2pt}\hrule height 0.3pt\vspace{3pt}
\textbf{Guide / Hint:} \\
{\small Initially $a=b=1$. Let statements be $T_1: [S_1: a=a+1, S_2: b=b+1]$ and $T_2: [S_3: b=2b, S_4: a=2a]$. Executing $T_1$ then $T_2$ yields $a=2, b=2 \to b=4, a=4 \implies (4, 4)$. Executing $T_2$ then $T_1$ yields $b=2, a=2 \to a=3, b=3 \implies (3, 3)$. Executing $S_1, S_3, S_4, S_2$ yields $a=2 \to b=2 \to a=4 \to b=3 \implies (4, 3)$. No valid interleaving can produce $(a=3, b=4)$.}

\end{tcolorbox}
\vspace{6pt}\hrule height 0.35pt\vspace{6pt}
\noindent
\begin{tcolorbox}[
    blanker,
    sidebyside,
    sidebyside align=top,
    lefthand width=0.55\textwidth,
    top=4pt, bottom=4pt, left=0pt, right=0pt,
    middle=16pt,
    underlay={
        \draw[dashed, black!80, line width=0.8pt, dash pattern=on 4pt off 3pt]
            ([xshift=0.55\textwidth+8pt]frame.north west) --
            ([xshift=0.55\textwidth+8pt]frame.south west);
    }
]
\textbf{Q.21} \quad {\footnotesize\textsc{[NAT $\cdot$ 2 Mark(s) $\cdot$ No Neg]}} \hfill {\footnotesize\textbf{GATE 2024 CS1} \quad \tiny\texttt{2024\_CS1\_CS\_Q57}}
\vspace{2pt}\hrule height 0.3pt\vspace{4pt}

Consider the following code snippet using the `fork()` and `wait()` system calls. Assume that the code compiles and runs correctly, and that the system calls run successfully without any errors.
\begin{lstlisting}[language=C]
int x = 3;
while(x > 0) {
    fork();
    printf("hello");
    wait(NULL);
    x--;
}
\end{lstlisting}
The total number of times the `printf` statement is executed is \_\_\_\_\_\_\_\. (Answer in integer)

\par\vspace{6pt}\noindent\textbf{Your Answer:} \fbox{\rule{0pt}{1.3ex}\hspace{3cm}}

\tcblower

\textbf{Official Key:} \textbf{\large 14 to 14} \hfill \textit{\small [Medium]}
\vspace{2pt}\hrule height 0.3pt\vspace{3pt}
\textbf{Topic:} Operating System \\
\textbf{Subtopic:} \textit{System calls, processes, threads, inter-process communication}
\vspace{2pt}\hrule height 0.3pt\vspace{3pt}
\textbf{Guide / Hint:} \\
{\small Initially 1 process with $x=3$. Iteration 1 ($x=3$): `fork()` creates 1 child (total 2 processes); both print "hello" (2 prints) and decrement $x$ to 2. Iteration 2 ($x=2$): each of the 2 processes forks (total 4 processes); all 4 print "hello" (4 prints) and decrement $x$ to 1. Iteration 3 ($x=1$): each of the 4 processes forks (total 8 processes); all 8 print "hello" (8 prints) and decrement $x$ to 0, whereupon the loop ends. Total prints $= 2 + 4 + 8 = 14$.}

\end{tcolorbox}
\vspace{6pt}\hrule height 0.35pt\vspace{6pt}
\noindent
\begin{tcolorbox}[
    blanker,
    sidebyside,
    sidebyside align=top,
    lefthand width=0.55\textwidth,
    top=4pt, bottom=4pt, left=0pt, right=0pt,
    middle=16pt,
    underlay={
        \draw[dashed, black!80, line width=0.8pt, dash pattern=on 4pt off 3pt]
            ([xshift=0.55\textwidth+8pt]frame.north west) --
            ([xshift=0.55\textwidth+8pt]frame.south west);
    }
]
\textbf{Q.22} \quad {\footnotesize\textsc{[NAT $\cdot$ 2 Mark(s) $\cdot$ No Neg]}} \hfill {\footnotesize\textbf{GATE 2024 CS1} \quad \tiny\texttt{2024\_CS1\_CS\_Q62}}
\vspace{2pt}\hrule height 0.3pt\vspace{4pt}

Consider a memory management system that uses a page size of 2 KB. Assume that both the physical and virtual addresses start from 0. Assume that the pages 0, 1, 2, and 3 are stored in the page frames 1, 3, 2, and 0, respectively. The physical address (in decimal format) corresponding to the virtual address 2500 (in decimal format) is \_\_\_\_\_\_\_\_\_\. (Answer in integer)

\par\vspace{6pt}\noindent\textbf{Your Answer:} \fbox{\rule{0pt}{1.3ex}\hspace{3cm}}

\tcblower

\textbf{Official Key:} \textbf{\large 6596 to 6596} \hfill \textit{\small [Easy]}
\vspace{2pt}\hrule height 0.3pt\vspace{3pt}
\textbf{Topic:} Operating System \\
\textbf{Subtopic:} \textit{Memory management and virtual memory}
\vspace{2pt}\hrule height 0.3pt\vspace{3pt}
\textbf{Guide / Hint:} \\
{\small With page size $2\text{ KB} = 2048\text{ bytes}$, virtual address $2500$ has virtual page number $\lfloor 2500 / 2048 \rfloor = 1$ and offset $2500 \pmod{2048} = 452$. Page 1 maps to frame 3. The physical address is $(\text{frame } 3 \times 2048) + \text{offset } 452 = 6144 + 452 = 6596$.}

\end{tcolorbox}
\vspace{6pt}\hrule height 0.35pt\vspace{6pt}
\noindent
\begin{tcolorbox}[
    blanker,
    sidebyside,
    sidebyside align=top,
    lefthand width=0.55\textwidth,
    top=4pt, bottom=4pt, left=0pt, right=0pt,
    middle=16pt,
    underlay={
        \draw[dashed, black!80, line width=0.8pt, dash pattern=on 4pt off 3pt]
            ([xshift=0.55\textwidth+8pt]frame.north west) --
            ([xshift=0.55\textwidth+8pt]frame.south west);
    }
]
\textbf{Q.23} \quad {\footnotesize\textsc{[MSQ $\cdot$ 1 Mark(s) $\cdot$ No Neg]}} \hfill {\footnotesize\textbf{GATE 2024 CS2} \quad \tiny\texttt{2024\_CS2\_CS\_Q24}}
\vspace{2pt}\hrule height 0.3pt\vspace{4pt}

Which of the following tasks is/are the responsibility/responsibilities of the memory management unit (MMU) in a system with paging-based memory management?

\par\vspace{6pt}\noindent
\begin{tabular}{@{}p{0.98\linewidth}@{}}
$\square$\ \textbf{(A)} Allocate a new page table for a newly created process \\
\addlinespace[3pt]
$\square$\ \textbf{(B)} Translate a virtual address to a physical address using the page table \\
\addlinespace[3pt]
$\square$\ \textbf{(C)} Raise a trap when a virtual address is not found in the page table \\
\addlinespace[3pt]
$\square$\ \textbf{(D)} Raise a trap when a process tries to write to a page marked with read-only permission in the page table
\end{tabular}

\tcblower

\textbf{Official Key:} \textbf{\large B;C;D} \hfill \textit{\small [Medium]}
\vspace{2pt}\hrule height 0.3pt\vspace{3pt}
\textbf{Topic:} Operating System \\
\textbf{Subtopic:} \textit{Memory management and virtual memory}
\vspace{2pt}\hrule height 0.3pt\vspace{3pt}
\textbf{Guide / Hint:} \\
{\small The MMU hardware translates virtual addresses to physical addresses using page tables/TLB (B), generates a page-fault trap when a page table entry is invalid (C), and raises a protection trap on access permission violations such as writing to read-only pages (D). Page table allocation (A) is an OS kernel duty.}

\end{tcolorbox}
\vspace{6pt}\hrule height 0.35pt\vspace{6pt}
\noindent
\begin{tcolorbox}[
    blanker,
    sidebyside,
    sidebyside align=top,
    lefthand width=0.55\textwidth,
    top=4pt, bottom=4pt, left=0pt, right=0pt,
    middle=16pt,
    underlay={
        \draw[dashed, black!80, line width=0.8pt, dash pattern=on 4pt off 3pt]
            ([xshift=0.55\textwidth+8pt]frame.north west) --
            ([xshift=0.55\textwidth+8pt]frame.south west);
    }
]
\textbf{Q.24} \quad {\footnotesize\textsc{[MSQ $\cdot$ 1 Mark(s) $\cdot$ No Neg]}} \hfill {\footnotesize\textbf{GATE 2024 CS2} \quad \tiny\texttt{2024\_CS2\_CS\_Q25}}
\vspace{2pt}\hrule height 0.3pt\vspace{4pt}

Consider a process P running on a CPU. Which one or more of the following events will always trigger a context switch by the OS that results in process P moving to a non-running state (e.g., ready, blocked)?

\par\vspace{6pt}\noindent
\begin{tabular}{@{}p{0.98\linewidth}@{}}
$\square$\ \textbf{(A)} P makes a blocking system call to read a block of data from the disk \\
\addlinespace[3pt]
$\square$\ \textbf{(B)} P tries to access a page that is in the swap space, triggering a page fault \\
\addlinespace[3pt]
$\square$\ \textbf{(C)} An interrupt is raised by the disk to deliver data requested by some other process \\
\addlinespace[3pt]
$\square$\ \textbf{(D)} A timer interrupt is raised by the hardware
\end{tabular}

\tcblower

\textbf{Official Key:} \textbf{\large A;B} \hfill \textit{\small [Medium]}
\vspace{2pt}\hrule height 0.3pt\vspace{3pt}
\textbf{Topic:} Operating System \\
\textbf{Subtopic:} \textit{System calls, processes, threads, inter-process communication}
\vspace{2pt}\hrule height 0.3pt\vspace{3pt}
\textbf{Guide / Hint:} \\
{\small When process P executes a blocking I/O system call, it must wait for disk transfer and is moved to the blocked state (A). When P experiences a page fault on a swapped page, it is put into the blocked state until disk I/O completes (B). Interrupts from other devices (C) or timer ticks (D) in preemptive systems might return directly to P after ISR handling if P has highest priority.}

\end{tcolorbox}
\vspace{6pt}\hrule height 0.35pt\vspace{6pt}
\noindent
\begin{tcolorbox}[
    blanker,
    sidebyside,
    sidebyside align=top,
    lefthand width=0.55\textwidth,
    top=4pt, bottom=4pt, left=0pt, right=0pt,
    middle=16pt,
    underlay={
        \draw[dashed, black!80, line width=0.8pt, dash pattern=on 4pt off 3pt]
            ([xshift=0.55\textwidth+8pt]frame.north west) --
            ([xshift=0.55\textwidth+8pt]frame.south west);
    }
]
\textbf{Q.25} \quad {\footnotesize\textsc{[MCQ $\cdot$ 2 Mark(s) $\cdot$ Neg: 0.66]}} \hfill {\footnotesize\textbf{GATE 2024 CS2} \quad \tiny\texttt{2024\_CS2\_CS\_Q37}}
\vspace{2pt}\hrule height 0.3pt\vspace{4pt}

Consider a single processor system with four processes A, B, C, and D, represented as given below, where for each process the first value is its arrival time, and the second value is its CPU burst time:
A (0, 10), B (2, 6), C (4, 3), and D (6, 7).
Which one of the following options gives the average waiting time for Shortest Remaining Time First (SRTF) and Non-Preemptive Shortest Job First (NP-SJF) scheduling?

\par\vspace{6pt}\noindent
\begin{tabular}{@{}p{0.48\linewidth}@{\hspace{0.04\linewidth}}p{0.48\linewidth}@{}}
$\square$\ \textbf{(A)} SRTF = 6, NP-SJF = 7 & $\square$\ \textbf{(B)} SRTF = 6, NP-SJF = 7.5 \\
\addlinespace[3pt]
$\square$\ \textbf{(C)} SRTF = 7, NP-SJF = 7.5 & $\square$\ \textbf{(D)} SRTF = 7, NP-SJF = 8.5
\end{tabular}

\tcblower

\textbf{Official Key:} \textbf{\large B} \hfill \textit{\small [Medium]}
\vspace{2pt}\hrule height 0.3pt\vspace{3pt}
\textbf{Topic:} Operating System \\
\textbf{Subtopic:} \textit{CPU and I/O scheduling}
\vspace{2pt}\hrule height 0.3pt\vspace{3pt}
\textbf{Guide / Hint:} \\
{\small Under SRTF: execution order is A(0-2), B(2-4), C(4-7), B(7-11), D(11-18), A(18-26). Turnaround times: A=26, B=9, C=3, D=12; Waiting times: A=16, B=3, C=0, D=5. Average waiting time $= 24/4 = 6.0$. Under NP-SJF: A runs non-preemptively (0-10), then C(10-13), B(13-19), D(19-26). Waiting times: A=0, C=6, B=11, D=13. Average waiting time $= 30/4 = 7.5$.}

\end{tcolorbox}
\vspace{6pt}\hrule height 0.35pt\vspace{6pt}
\noindent
\begin{tcolorbox}[
    blanker,
    sidebyside,
    sidebyside align=top,
    lefthand width=0.55\textwidth,
    top=4pt, bottom=4pt, left=0pt, right=0pt,
    middle=16pt,
    underlay={
        \draw[dashed, black!80, line width=0.8pt, dash pattern=on 4pt off 3pt]
            ([xshift=0.55\textwidth+8pt]frame.north west) --
            ([xshift=0.55\textwidth+8pt]frame.south west);
    }
]
\textbf{Q.26} \quad {\footnotesize\textsc{[MSQ $\cdot$ 2 Mark(s) $\cdot$ No Neg]}} \hfill {\footnotesize\textbf{GATE 2024 CS2} \quad \tiny\texttt{2024\_CS2\_CS\_Q46}}
\vspace{2pt}\hrule height 0.3pt\vspace{4pt}

Consider a multi-threaded program with two threads T1 and T2. The threads share two semaphores: s1 (initialized to 1) and s2 (initialized to 0). The threads also share a global variable x (initialized to 0). The threads execute the code shown below:
\begin{verbatim}
// code of T1        // code of T2
wait(s1);            wait(s1);
x = x + 1;           x = x + 1;
print(x);            print(x);
wait(s2);            signal(s2);
signal(s1);          signal(s1);
\end{verbatim}
Which of the following outcomes is/are possible when threads T1 and T2 execute concurrently?

\par\vspace{6pt}\noindent
\begin{tabular}{@{}p{0.98\linewidth}@{}}
$\square$\ \textbf{(A)} T1 runs first and prints 1, T2 runs next and prints 2 \\
\addlinespace[3pt]
$\square$\ \textbf{(B)} T2 runs first and prints 1, T1 runs next and prints 2 \\
\addlinespace[3pt]
$\square$\ \textbf{(C)} T1 runs first and prints 1, T2 does not print anything (deadlock) \\
\addlinespace[3pt]
$\square$\ \textbf{(D)} T2 runs first and prints 1, T1 does not print anything (deadlock)
\end{tabular}

\tcblower

\textbf{Official Key:} \textbf{\large B;C} \hfill \textit{\small [Medium]}
\vspace{2pt}\hrule height 0.3pt\vspace{3pt}
\textbf{Topic:} Operating System \\
\textbf{Subtopic:} \textit{Concurrency and synchronization}
\vspace{2pt}\hrule height 0.3pt\vspace{3pt}
\textbf{Guide / Hint:} \\
{\small If T2 runs first: T2 takes s1, increments x (x=1), prints 1, signals s2, signals s1. Then T1 executes: takes s1, increments x (x=2), prints 2, completes wait(s2) without blocking (since s2 was signaled), signals s1 (outcome B). If T1 runs first: T1 takes s1, prints 1, and blocks on wait(s2) while holding s1, causing deadlock where T2 blocks on wait(s1) (outcome C).}

\end{tcolorbox}
\vspace{6pt}\hrule height 0.35pt\vspace{6pt}
\noindent
\begin{tcolorbox}[
    blanker,
    sidebyside,
    sidebyside align=top,
    lefthand width=0.55\textwidth,
    top=4pt, bottom=4pt, left=0pt, right=0pt,
    middle=16pt,
    underlay={
        \draw[dashed, black!80, line width=0.8pt, dash pattern=on 4pt off 3pt]
            ([xshift=0.55\textwidth+8pt]frame.north west) --
            ([xshift=0.55\textwidth+8pt]frame.south west);
    }
]
\textbf{Q.27} \quad {\footnotesize\textsc{[NAT $\cdot$ 2 Mark(s) $\cdot$ No Neg]}} \hfill {\footnotesize\textbf{GATE 2024 CS2} \quad \tiny\texttt{2024\_CS2\_CS\_Q53}}
\vspace{2pt}\hrule height 0.3pt\vspace{4pt}

Consider a disk with the following specifications: rotation speed of 6000 RPM, average seek time of 5 milliseconds, 500 sectors/track, 512-byte sectors. A file has content stored in 3000 sectors located randomly on the disk. Assuming average rotational latency, the total time (in seconds, rounded off to 2 decimal places) to read the entire file from the disk is \_\_\_\_\_\_\. (Answer in range 29.50 to 30.50)

\par\vspace{6pt}\noindent\textbf{Your Answer:} \fbox{\rule{0pt}{1.3ex}\hspace{3cm}}

\tcblower

\textbf{Official Key:} \textbf{\large 29.50 to 30.50} \hfill \textit{\small [Medium]}
\vspace{2pt}\hrule height 0.3pt\vspace{3pt}
\textbf{Topic:} Operating System \\
\textbf{Subtopic:} \textit{File systems}
\vspace{2pt}\hrule height 0.3pt\vspace{3pt}
\textbf{Guide / Hint:} \\
{\small At 6000 RPM, one rotation takes $\frac{60}{6000} = 10\text{ ms}$. Average rotational latency $= 10/2 = 5\text{ ms}$. Transfer time per sector $= \frac{10\text{ ms}}{500} = 0.02\text{ ms}$. Time to access each random sector $= \text{seek} + \text{rotational latency} + \text{transfer} = 5 + 5 + 0.02 = 10.02\text{ ms}$. Total time for 3000 random sectors $= 3000 \times 10.02\text{ ms} = 30060\text{ ms} \approx 30.06$ seconds.}

\end{tcolorbox}
\vspace{6pt}\hrule height 0.35pt\vspace{6pt}
\noindent
\begin{tcolorbox}[
    blanker,
    sidebyside,
    sidebyside align=top,
    lefthand width=0.55\textwidth,
    top=4pt, bottom=4pt, left=0pt, right=0pt,
    middle=16pt,
    underlay={
        \draw[dashed, black!80, line width=0.8pt, dash pattern=on 4pt off 3pt]
            ([xshift=0.55\textwidth+8pt]frame.north west) --
            ([xshift=0.55\textwidth+8pt]frame.south west);
    }
]
\textbf{Q.28} \quad {\footnotesize\textsc{[NAT $\cdot$ 2 Mark(s) $\cdot$ No Neg]}} \hfill {\footnotesize\textbf{GATE 2024 CS2} \quad \tiny\texttt{2024\_CS2\_CS\_Q64}}
\vspace{2pt}\hrule height 0.3pt\vspace{4pt}

Consider a 32-bit system with 4 KB page size and page table entries of size 4 bytes each. Assume $1\text{ KB} = 2^{10}$ bytes. The OS uses a 2-level page table for memory management, with an outer page directory and an inner page table. The OS allocates 1 page for the outer page directory upon process creation. The OS uses demand paging when allocating memory for the inner page table. An active process accesses 2000 unique pages during its execution. Let $X$ denote the minimum and $Y$ denote the maximum number of pages across the two levels of the page table of the process. The value of $X + Y$ is \_\_\_\_\_\_\. (Answer in integer)

\par\vspace{6pt}\noindent\textbf{Your Answer:} \fbox{\rule{0pt}{1.3ex}\hspace{3cm}}

\tcblower

\textbf{Official Key:} \textbf{\large 1028 to 1028} \hfill \textit{\small [Hard]}
\vspace{2pt}\hrule height 0.3pt\vspace{3pt}
\textbf{Topic:} Operating System \\
\textbf{Subtopic:} \textit{Memory management and virtual memory}
\vspace{2pt}\hrule height 0.3pt\vspace{3pt}
\textbf{Guide / Hint:} \\
{\small Page size is $4\text{ KB} = 2^{12}$ bytes, PTE size is 4 bytes, so each page table holds $2^{12} / 4 = 1024$ entries. The process accesses 2000 unique pages. Minimum inner page tables $X$: $\lceil 2000 / 1024 \rceil = 2$ inner pages plus 1 outer page directory $= 3$. Maximum inner page tables $Y$: 2000 pages distributed across 2000 different inner tables, bounded by total possible entries in directory $= 1024$ inner pages plus 1 outer page directory $= 1025$. Hence, $X + Y = 3 + 1025 = 1028$.}

\end{tcolorbox}
\vspace{6pt}\hrule height 0.35pt\vspace{6pt}
\noindent
\begin{tcolorbox}[
    blanker,
    sidebyside,
    sidebyside align=top,
    lefthand width=0.55\textwidth,
    top=4pt, bottom=4pt, left=0pt, right=0pt,
    middle=16pt,
    underlay={
        \draw[dashed, black!80, line width=0.8pt, dash pattern=on 4pt off 3pt]
            ([xshift=0.55\textwidth+8pt]frame.north west) --
            ([xshift=0.55\textwidth+8pt]frame.south west);
    }
]
\textbf{Q.29} \quad {\footnotesize\textsc{[MCQ $\cdot$ 1 Mark(s) $\cdot$ Neg: 0.33]}} \hfill {\footnotesize\textbf{GATE 2025 CS1} \quad \tiny\texttt{2025\_CS1\_CS\_Q14}}
\vspace{2pt}\hrule height 0.3pt\vspace{4pt}

Consider a demand paging memory management system with 32-bit logical address, 20-bit physical address, and page size of 2048 bytes. Assuming that the memory is byte addressable, what is the maximum number of entries in the page table?

\par\vspace{6pt}\noindent
\begin{tabular}{@{}p{0.48\linewidth}@{\hspace{0.04\linewidth}}p{0.48\linewidth}@{}}
$\square$\ \textbf{(A)} $2^{21}$ & $\square$\ \textbf{(B)} $2^{20}$ \\
\addlinespace[3pt]
$\square$\ \textbf{(C)} $2^{22}$ & $\square$\ \textbf{(D)} $2^{24}$
\end{tabular}

\tcblower

\textbf{Official Key:} \textbf{\large A} \hfill \textit{\small [Medium]}
\vspace{2pt}\hrule height 0.3pt\vspace{3pt}
\textbf{Topic:} Operating System \\
\textbf{Subtopic:} \textit{Memory management and virtual memory}
\vspace{2pt}\hrule height 0.3pt\vspace{3pt}
\textbf{Guide / Hint:} \\
{\small With a page size of $2048 = 2^{11}$ bytes, the page offset requires 11 bits. The 32-bit logical address is split into page number and page offset, so the page number has $32 - 11 = 21$ bits. Hence, the page table contains at most $2^{21}$ entries.}

\end{tcolorbox}
\vspace{6pt}\hrule height 0.35pt\vspace{6pt}
\noindent
\begin{tcolorbox}[
    blanker,
    sidebyside,
    sidebyside align=top,
    lefthand width=0.55\textwidth,
    top=4pt, bottom=4pt, left=0pt, right=0pt,
    middle=16pt,
    underlay={
        \draw[dashed, black!80, line width=0.8pt, dash pattern=on 4pt off 3pt]
            ([xshift=0.55\textwidth+8pt]frame.north west) --
            ([xshift=0.55\textwidth+8pt]frame.south west);
    }
]
\textbf{Q.30} \quad {\footnotesize\textsc{[NAT $\cdot$ 1 Mark(s) $\cdot$ No Neg]}} \hfill {\footnotesize\textbf{GATE 2025 CS1} \quad \tiny\texttt{2025\_CS1\_CS\_Q29}}
\vspace{2pt}\hrule height 0.3pt\vspace{4pt}

Suppose in a multiprogramming environment, the following C program segment is executed. A process goes into I/O queue whenever an I/O related operation is performed. Assume that there will always be a context switch whenever a process requests for an I/O, and also whenever the process returns from an I/O. The number of times the process will enter the ready queue during its lifetime (not counting the time the process enters the ready queue when it is run initially) is \_\_\_\_\_\_\. (Answer in integer)
\par\medskip
\begin{verbatim}
int main()
{
    int x=0, i=0;
    scanf("%d", &x);
    for(i=0; i<20; i++)
    {
        x = x + 20;
        printf("%d\n", x);
    }
    return 0;
}
\end{verbatim}

\par\vspace{6pt}\noindent\textbf{Your Answer:} \fbox{\rule{0pt}{1.3ex}\hspace{3cm}}

\tcblower

\textbf{Official Key:} \textbf{\large 21 to 21} \hfill \textit{\small [Medium]}
\vspace{2pt}\hrule height 0.3pt\vspace{3pt}
\textbf{Topic:} Operating System \\
\textbf{Subtopic:} \textit{CPU and I/O scheduling}
\vspace{2pt}\hrule height 0.3pt\vspace{3pt}
\textbf{Guide / Hint:} \\
{\small There is 1 initial `scanf` and 20 iterations of `printf`, totaling $1 + 20 = 21$ I/O operations. Each time an I/O operation completes, the process transitions from the I/O (waiting) queue back into the ready queue. Excluding the initial program launch into the ready queue, the process enters the ready queue exactly 21 times.}

\end{tcolorbox}
\vspace{6pt}\hrule height 0.35pt\vspace{6pt}
\noindent
\begin{tcolorbox}[
    blanker,
    sidebyside,
    sidebyside align=top,
    lefthand width=0.55\textwidth,
    top=4pt, bottom=4pt, left=0pt, right=0pt,
    middle=16pt,
    underlay={
        \draw[dashed, black!80, line width=0.8pt, dash pattern=on 4pt off 3pt]
            ([xshift=0.55\textwidth+8pt]frame.north west) --
            ([xshift=0.55\textwidth+8pt]frame.south west);
    }
]
\textbf{Q.31} \quad {\footnotesize\textsc{[MCQ $\cdot$ 2 Mark(s) $\cdot$ Neg: 0.66]}} \hfill {\footnotesize\textbf{GATE 2025 CS1} \quad \tiny\texttt{2025\_CS1\_CS\_Q38}}
\vspace{2pt}\hrule height 0.3pt\vspace{4pt}

A computer has two processors, $M_1$ and $M_2$. Four processes $P_1, P_2, P_3, P_4$ with CPU bursts of 20, 16, 25, and 10 milliseconds, respectively, arrive at the same time and these are the only processes in the system. The scheduler uses non-preemptive priority scheduling, with priorities decided as follows:
\begin{itemize}
\item $M_1$ uses priority of execution for the processes as, $P_1 > P_3 > P_2 > P_4$, i.e., $P_1$ and $P_4$ have highest and lowest priorities, respectively.
\item $M_2$ uses priority of execution for the processes as, $P_2 > P_3 > P_4 > P_1$, i.e., $P_2$ and $P_1$ have highest and lowest priorities, respectively.
\end{itemize}
A process $P_i$ is scheduled to a processor $M_k$, if the processor is free and no other process $P_j$ is waiting with higher priority. At any given point of time, a process can be allocated to any one of the free processors without violating the execution priority rules. Ignore the context switch time. What will be the average waiting time of the processes in milliseconds?

\par\vspace{6pt}\noindent
\begin{tabular}{@{}p{0.48\linewidth}@{\hspace{0.04\linewidth}}p{0.48\linewidth}@{}}
$\square$\ \textbf{(A)} 9.00 & $\square$\ \textbf{(B)} 8.75 \\
\addlinespace[3pt]
$\square$\ \textbf{(C)} 6.50 & $\square$\ \textbf{(D)} 7.50
\end{tabular}

\tcblower

\textbf{Official Key:} \textbf{\large A} \hfill \textit{\small [Hard]}
\vspace{2pt}\hrule height 0.3pt\vspace{3pt}
\textbf{Topic:} Operating System \\
\textbf{Subtopic:} \textit{CPU and I/O scheduling}
\vspace{2pt}\hrule height 0.3pt\vspace{3pt}
\textbf{Guide / Hint:} \\
{\small At $t = 0$, $M_1$ takes highest priority process $P_1$ (burst 20) and $M_2$ takes $P_2$ (burst 16). At $t = 16$, $M_2$ frees up; among remaining processes $\{P_3, P_4\}$, $P_3$ has higher priority on $M_2$, so $P_3$ runs on $M_2$ ($t = 16$ to 41). At $t = 20$, $M_1$ frees up and executes $P_4$ ($t = 20$ to 30). Waiting times are: $P_1 = 0$, $P_2 = 0$, $P_3 = 16$, $P_4 = 20$. Average waiting time is $(0 + 0 + 16 + 20) / 4 = 36 / 4 = 9.00$ ms.}

\end{tcolorbox}
\vspace{6pt}\hrule height 0.35pt\vspace{6pt}
\noindent
\begin{tcolorbox}[
    blanker,
    sidebyside,
    sidebyside align=top,
    lefthand width=0.55\textwidth,
    top=4pt, bottom=4pt, left=0pt, right=0pt,
    middle=16pt,
    underlay={
        \draw[dashed, black!80, line width=0.8pt, dash pattern=on 4pt off 3pt]
            ([xshift=0.55\textwidth+8pt]frame.north west) --
            ([xshift=0.55\textwidth+8pt]frame.south west);
    }
]
\textbf{Q.32} \quad {\footnotesize\textsc{[NAT $\cdot$ 2 Mark(s) $\cdot$ No Neg]}} \hfill {\footnotesize\textbf{GATE 2025 CS1} \quad \tiny\texttt{2025\_CS1\_CS\_Q51}}
\vspace{2pt}\hrule height 0.3pt\vspace{4pt}

A disk of size 512M bytes is divided into blocks of 64K bytes. A file is stored in the disk using linked allocation. In linked allocation, each data block reserves 4 bytes to store the pointer to the next data block. The link part of the last data block contains a NULL pointer (also of 4 bytes). Suppose a file of 1M bytes needs to be stored in the disk. Assume $1\text{K} = 2^{10}$ and $1\text{M} = 2^{20}$. The amount of space in bytes that will be wasted due to internal fragmentation is \_\_\_\_\_\_\. (Answer in integer)

\par\vspace{6pt}\noindent\textbf{Your Answer:} \fbox{\rule{0pt}{1.3ex}\hspace{3cm}}

\tcblower

\textbf{Official Key:} \textbf{\large 65468 to 65468} \hfill \textit{\small [Medium]}
\vspace{2pt}\hrule height 0.3pt\vspace{3pt}
\textbf{Topic:} Operating System \\
\textbf{Subtopic:} \textit{File systems}
\vspace{2pt}\hrule height 0.3pt\vspace{3pt}
\textbf{Guide / Hint:} \\
{\small Block size is $64\text{ KB} = 65536$ bytes, with 4 bytes for pointer, leaving $65536 - 4 = 65532$ bytes of file data per block. File size is $1\text{ MB} = 1048576$ bytes. Number of blocks needed is $\lceil 1048576 / 65532 \rceil = \lceil 16.001037 \rceil = 17$ blocks. The first 16 blocks store $16 \times 65532 = 1048512$ bytes, leaving $1048576 - 1048512 = 64$ bytes stored in the 17th block. Internal fragmentation in the last block is $65532 - 64 = 65468$ bytes.}

\end{tcolorbox}
\vspace{6pt}\hrule height 0.35pt\vspace{6pt}
\noindent
\begin{tcolorbox}[
    blanker,
    sidebyside,
    sidebyside align=top,
    lefthand width=0.55\textwidth,
    top=4pt, bottom=4pt, left=0pt, right=0pt,
    middle=16pt,
    underlay={
        \draw[dashed, black!80, line width=0.8pt, dash pattern=on 4pt off 3pt]
            ([xshift=0.55\textwidth+8pt]frame.north west) --
            ([xshift=0.55\textwidth+8pt]frame.south west);
    }
]
\textbf{Q.33} \quad {\footnotesize\textsc{[NAT $\cdot$ 2 Mark(s) $\cdot$ No Neg]}} \hfill {\footnotesize\textbf{GATE 2025 CS1} \quad \tiny\texttt{2025\_CS1\_CS\_Q54}}
\vspace{2pt}\hrule height 0.3pt\vspace{4pt}

In optimal page replacement algorithm, information about all future page references is available to the operating system (OS). A modification of the optimal page replacement algorithm is as follows:
The OS correctly predicts only up to next 4 page references (including the current page) at the time of allocating a frame to a page.
A process accesses the pages in the following order of page numbers:
1, 3, 2, 4, 2, 3, 1, 2, 4, 3, 1, 4.
If the system has three memory frames that are initially empty, the number of page faults that will occur during execution of the process is \_\_\_\_\_\_\. (Answer in integer)

\par\vspace{6pt}\noindent\textbf{Your Answer:} \fbox{\rule{0pt}{1.3ex}\hspace{3cm}}

\tcblower

\textbf{Official Key:} \textbf{\large 6 to 6} \hfill \textit{\small [Medium]}
\vspace{2pt}\hrule height 0.3pt\vspace{3pt}
\textbf{Topic:} Operating System \\
\textbf{Subtopic:} \textit{Memory management and virtual memory}
\vspace{2pt}\hrule height 0.3pt\vspace{3pt}
\textbf{Guide / Hint:} \\
{\small Simulating the 3 frames with 4-step lookahead: 1 (fault, [1]), 3 (fault, [1, 3]), 2 (fault, [1, 3, 2]). At page 4 (fault): window is (4, 2, 3, 1), farthest is 1, so replace 1 -> [4, 3, 2]. Pages 2, 3 hit. At page 1 (fault): window (1, 2, 4, 3), replace 3 -> [4, 1, 2]. Page 2, 4, 3 (fault): replace 1 -> [4, 3, 2]. Total faults = 3 compulsory + 3 replacement = 6 page faults.}

\end{tcolorbox}
\vspace{6pt}\hrule height 0.35pt\vspace{6pt}
\noindent
\begin{tcolorbox}[
    blanker,
    sidebyside,
    sidebyside align=top,
    lefthand width=0.55\textwidth,
    top=4pt, bottom=4pt, left=0pt, right=0pt,
    middle=16pt,
    underlay={
        \draw[dashed, black!80, line width=0.8pt, dash pattern=on 4pt off 3pt]
            ([xshift=0.55\textwidth+8pt]frame.north west) --
            ([xshift=0.55\textwidth+8pt]frame.south west);
    }
]
\textbf{Q.34} \quad {\footnotesize\textsc{[MCQ $\cdot$ 1 Mark(s) $\cdot$ Neg: 0.33]}} \hfill {\footnotesize\textbf{GATE 2025 CS2} \quad \tiny\texttt{2025\_CS2\_CS\_Q26}}
\vspace{2pt}\hrule height 0.3pt\vspace{4pt}

Processes $P_1, P_2, P_3, P_4$ arrive in that order at times 0, 1, 2, and 8 milliseconds respectively, and have execution times of 10, 13, 6, and 9 milliseconds respectively. Shortest Remaining Time First (SRTF) algorithm is used as the CPU scheduling policy. Ignore context switching times. Which ONE of the following correctly gives the average turnaround time of the four processes in milliseconds?

\par\vspace{6pt}\noindent
\begin{tabular}{@{}p{0.48\linewidth}@{\hspace{0.04\linewidth}}p{0.48\linewidth}@{}}
$\square$\ \textbf{(A)} 22 & $\square$\ \textbf{(B)} 15 \\
\addlinespace[3pt]
$\square$\ \textbf{(C)} 37 & $\square$\ \textbf{(D)} 19
\end{tabular}

\tcblower

\textbf{Official Key:} \textbf{\large D} \hfill \textit{\small [Medium]}
\vspace{2pt}\hrule height 0.3pt\vspace{3pt}
\textbf{Topic:} Operating System \\
\textbf{Subtopic:} \textit{CPU and I/O scheduling: SRTF}
\vspace{2pt}\hrule height 0.3pt\vspace{3pt}
\textbf{Guide / Hint:} \\
{\small Using preemptive SRTF: $P_1$ runs [0,2], $P_3$ runs [2,8] (finishes at 8, TAT = $8-2=6$). Then $P_4$ (remaining 9) runs [8,17] (finishes at 17, TAT = $17-8=9$). $P_1$ (remaining 8) runs [17,25] (finishes at 25, TAT = $25-0=25$). $P_2$ runs [25,38] (finishes at 38, TAT = $38-1=37$). Average TAT = $\frac{25 + 37 + 6 + 9}{4} = \frac{76}{4} = 19$ ms.}

\end{tcolorbox}
\vspace{6pt}\hrule height 0.35pt\vspace{6pt}
\noindent
\begin{tcolorbox}[
    blanker,
    sidebyside,
    sidebyside align=top,
    lefthand width=0.55\textwidth,
    top=4pt, bottom=4pt, left=0pt, right=0pt,
    middle=16pt,
    underlay={
        \draw[dashed, black!80, line width=0.8pt, dash pattern=on 4pt off 3pt]
            ([xshift=0.55\textwidth+8pt]frame.north west) --
            ([xshift=0.55\textwidth+8pt]frame.south west);
    }
]
\textbf{Q.35} \quad {\footnotesize\textsc{[MSQ $\cdot$ 2 Mark(s) $\cdot$ No Neg]}} \hfill {\footnotesize\textbf{GATE 2025 CS2} \quad \tiny\texttt{2025\_CS2\_CS\_Q47}}
\vspace{2pt}\hrule height 0.3pt\vspace{4pt}

Consider a demand paging system with three frames, and the following page reference string: 1 2 3 4 5 4 1 6 4 2 1. Which page replacement policy produces the minimum number of page faults for this reference string?

\par\vspace{6pt}\noindent
\begin{tabular}{@{}p{0.98\linewidth}@{}}
$\square$\ \textbf{(A)} Least Recently Used page replacement policy \\
\addlinespace[3pt]
$\square$\ \textbf{(B)} Least Frequently Used page replacement policy \\
\addlinespace[3pt]
$\square$\ \textbf{(C)} Most Frequently Used page replacement policy \\
\addlinespace[3pt]
$\square$\ \textbf{(D)} Optimal page replacement policy
\end{tabular}

\tcblower

\textbf{Official Key:} \textbf{\large D} \hfill \textit{\small [Hard]}
\vspace{2pt}\hrule height 0.3pt\vspace{3pt}
\textbf{Topic:} Operating System \\
\textbf{Subtopic:} \textit{Memory management and virtual memory: page replacement algorithms}
\vspace{2pt}\hrule height 0.3pt\vspace{3pt}
\textbf{Guide / Hint:} \\
{\small Trace the page reference string through 3 frames under each algorithm. Optimal page replacement replaces the page that will not be used for the longest period of time in the future, matching the minimum hit/fault sequence.}

\end{tcolorbox}
\vspace{6pt}\hrule height 0.35pt\vspace{6pt}
\noindent
\begin{tcolorbox}[
    blanker,
    sidebyside,
    sidebyside align=top,
    lefthand width=0.55\textwidth,
    top=4pt, bottom=4pt, left=0pt, right=0pt,
    middle=16pt,
    underlay={
        \draw[dashed, black!80, line width=0.8pt, dash pattern=on 4pt off 3pt]
            ([xshift=0.55\textwidth+8pt]frame.north west) --
            ([xshift=0.55\textwidth+8pt]frame.south west);
    }
]
\textbf{Q.36} \quad {\footnotesize\textsc{[MSQ $\cdot$ 2 Mark(s) $\cdot$ No Neg]}} \hfill {\footnotesize\textbf{GATE 2025 CS2} \quad \tiny\texttt{2025\_CS2\_CS\_Q48}}
\vspace{2pt}\hrule height 0.3pt\vspace{4pt}

Let $P = \{P_1, P_2, P_3, P_4\}$ be active processes and $R = \{R_1, R_2, R_3, R_4\}$ be resources in an operating system. Given the allocation and request graph, which of the following actions makes the system deadlock free?

\par\vspace{6pt}\noindent
\begin{tabular}{@{}p{0.98\linewidth}@{}}
$\square$\ \textbf{(A)} Aborting $P_1$ makes the system deadlock free. \\
\addlinespace[3pt]
$\square$\ \textbf{(B)} Aborting $P_3$ makes the system deadlock free. \\
\addlinespace[3pt]
$\square$\ \textbf{(C)} Aborting $P_2$ makes the system deadlock free. \\
\addlinespace[3pt]
$\square$\ \textbf{(D)} Aborting $P_1$ and $P_4$ makes the system deadlock free.
\end{tabular}

\tcblower

\textbf{Official Key:} \textbf{\large C;D} \hfill \textit{\small [Hard]}
\vspace{2pt}\hrule height 0.3pt\vspace{3pt}
\textbf{Topic:} Operating System \\
\textbf{Subtopic:} \textit{Deadlock: Resource Allocation Graph and deadlock recovery}
\vspace{2pt}\hrule height 0.3pt\vspace{3pt}
\textbf{Guide / Hint:} \\
{\small In the Resource Allocation Graph, identify the directed cycles creating deadlock. Aborting process $P_2$ eliminates the dependency cycles directly, as does aborting both $P_1$ and $P_4$.}

\end{tcolorbox}
\vspace{6pt}\hrule height 0.35pt\vspace{6pt}
\noindent
\begin{tcolorbox}[
    blanker,
    sidebyside,
    sidebyside align=top,
    lefthand width=0.55\textwidth,
    top=4pt, bottom=4pt, left=0pt, right=0pt,
    middle=16pt,
    underlay={
        \draw[dashed, black!80, line width=0.8pt, dash pattern=on 4pt off 3pt]
            ([xshift=0.55\textwidth+8pt]frame.north west) --
            ([xshift=0.55\textwidth+8pt]frame.south west);
    }
]
\textbf{Q.37} \quad {\footnotesize\textsc{[NAT $\cdot$ 2 Mark(s) $\cdot$ No Neg]}} \hfill {\footnotesize\textbf{GATE 2025 CS2} \quad \tiny\texttt{2025\_CS2\_CS\_Q58}}
\vspace{2pt}\hrule height 0.3pt\vspace{4pt}

A computer system supports a logical address space of $2^{32}$ bytes. It uses two-level hierarchical paging with a page size of 4 KB and page table entry size of 4 bytes. The number of bits in the outer page table index is \_\_\_\_\_\_\_\_.

\par\vspace{6pt}\noindent\textbf{Your Answer:} \fbox{\rule{0pt}{1.3ex}\hspace{3cm}}

\tcblower

\textbf{Official Key:} \textbf{\large 11 to 11} \hfill \textit{\small [Hard]}
\vspace{2pt}\hrule height 0.3pt\vspace{3pt}
\textbf{Topic:} Operating System \\
\textbf{Subtopic:} \textit{Memory management and virtual memory: hierarchical paging}
\vspace{2pt}\hrule height 0.3pt\vspace{3pt}
\textbf{Guide / Hint:} \\
{\small For a $2^{32}$-byte logical address space with two-level hierarchical paging, page offset and page table index bits are derived from page size and entry size, yielding 11 bits for the first-level directory index.}

\end{tcolorbox}
\vspace{6pt}\hrule height 0.35pt\vspace{6pt}
\end{document}